\chapter{MVC + Service Layer Architecture}

\section{Professional Backend Folder Structure}

\begin{diagrambox}[backend/src Folder Tree]
\begin{verbatim}
backend/src/
|-- config/
|   `-- db.js               MongoDB connection setup
|-- controllers/
|   |-- auth.controller.js   HTTP layer: parses req, calls service, sends res
|   `-- task.controller.js
|-- middleware/
|   |-- auth.middleware.js   protect(), role checks
|   `-- error.middleware.js  centralized error handler
|-- models/
|   |-- User.js              Mongoose schema/model
|   `-- Task.js
|-- routes/
|   |-- auth.routes.js       path -> controller wiring
|   `-- task.routes.js
|-- services/
|   |-- auth.service.js      business logic + Mongoose queries
|   `-- task.service.js
|-- validators/
|   `-- task.validator.js    input validation schemas
|-- utils/
|   |-- generateToken.js
|   `-- asyncHandler.js
|-- app.js                   Express app + middleware + routes
`-- server.js                HTTP listener + DB connection
\end{verbatim}
\end{diagrambox}

\section{The Request Flow}

\begin{diagrambox}[Layered Request Flow]
\begin{verbatim}
Client Request
     |
     v
  Route          (path matching: POST /api/tasks)
     |
     v
 Middleware       (auth check, validation)
     |
     v
 Controller       (parse req, call service, shape res)
     |
     v
  Service         (business rules, orchestrates model calls)
     |
     v
   Model          (Mongoose schema methods: Task.create, Task.find)
     |
     v
 MongoDB           (actual document storage)
\end{verbatim}
\end{diagrambox}

\section{Anti-Pattern: Everything in the Route Handler}

\begin{warnbox}[ANTI-PATTERN]
The following is what an unstructured, "just make it work" route looks like. It mixes validation, auth, database access, and business logic in one function. This does not scale past a handful of endpoints.
\end{warnbox}

\begin{lstlisting}[language=JavaScript]
// routes/task.routes.js  -- DO NOT DO THIS
router.post("/", async (req, res) => {
  // auth check inline
  const token = req.cookies.token;
  if (!token) return res.status(401).json({ message: "Not authenticated" });
  let userId;
  try {
    const decoded = jwt.verify(token, process.env.JWT_SECRET);
    userId = decoded.id;
  } catch (e) {
    return res.status(401).json({ message: "Invalid token" });
  }

  // validation inline
  const { title, description, status, priority, dueDate } = req.body;
  if (!title || title.trim() === "") {
    return res.status(400).json({ message: "Title is required" });
  }
  if (status && !["todo", "progress", "completed"].includes(status)) {
    return res.status(400).json({ message: "Invalid status" });
  }

  // business logic + db access inline
  try {
    const existingCount = await Task.countDocuments({ user: userId });
    if (existingCount >= 100) {
      return res.status(400).json({ message: "Task limit reached" });
    }
    const task = new Task({ title, description, status, priority, dueDate, user: userId });
    await task.save();
    const populated = await Task.findById(task._id).populate("user", "name email");
    return res.status(201).json(populated);
  } catch (err) {
    return res.status(500).json({ message: err.message });
  }
});
\end{lstlisting}

\section{Refactored: Route -> Controller -> Service}

\subsection{routes/task.routes.js}

\begin{lstlisting}[language=JavaScript]
const express = require("express");
const { createTask } = require("../controllers/task.controller");
const { protect } = require("../middleware/auth.middleware");
const { validateTask } = require("../validators/task.validator");

const router = express.Router();

router.post("/", protect, validateTask, createTask);

module.exports = router;
\end{lstlisting}

\subsection{controllers/task.controller.js}

\begin{lstlisting}[language=JavaScript]
const taskService = require("../services/task.service");

// Controller: translates HTTP <-> service calls. No Mongoose here.
const createTask = async (req, res, next) => {
  try {
    const task = await taskService.createTask(req.body, req.user.id);
    res.status(201).json(task);
  } catch (err) {
    next(err); // delegate to centralized error middleware
  }
};

module.exports = { createTask };
\end{lstlisting}

\subsection{services/task.service.js}

\begin{lstlisting}[language=JavaScript]
const Task = require("../models/Task");

const MAX_TASKS_PER_USER = 100;

// Service: owns business rules and all Mongoose interaction.
const createTask = async (taskData, userId) => {
  const existingCount = await Task.countDocuments({ user: userId });
  if (existingCount >= MAX_TASKS_PER_USER) {
    const err = new Error("Task limit reached");
    err.statusCode = 400;
    throw err;
  }

  const task = await Task.create({ ...taskData, user: userId });
  return task.populate("user", "name email");
};

module.exports = { createTask };
\end{lstlisting}

\section{Layer Responsibilities}

\begin{table}[h]
\centering
\begin{tabularx}{\textwidth}{l X X}
\toprule
\textbf{Layer} & \textbf{Responsibility} & \textbf{Should NOT contain} \\
\midrule
Route       & Map HTTP method+path to a controller; attach middleware & Business logic, database calls, validation logic \\
Middleware  & Cross-cutting concerns: auth, validation, logging, error handling & Resource-specific business rules \\
Controller  & Read \texttt{req}, call one service function, shape \texttt{res} & Direct Mongoose queries, complex business rules \\
Service     & Business rules, orchestration, calls to one or more models & \texttt{req}/\texttt{res} objects, HTTP status codes \\
Model       & Schema definition, field-level validation, Mongoose hooks & Business rules that span multiple resources \\
\bottomrule
\end{tabularx}
\end{table}

\begin{tipbox}[TIP]
A good litmus test: if you could reuse the function from a CLI script or a cron job (no \texttt{req}/\texttt{res} in sight), it belongs in the service layer.
\end{tipbox}
